% Figura F5-2 del anexo CRISP-DM, con la leyenda adaptada a la memoria.
\begin{figure}[!htbp]
\centering
\pgfplotstableread[col sep=semicolon]{\graficasdir/F5-2.dat}\tablaG
\begin{tikzpicture}
\begin{axis}[width=0.93\linewidth, height=6.4cm, xlabel={Umbral $\tau$ sobre el puntaje del reordenador}, ylabel={Proporción}, ymin=0, ymax=1.05, grid=major, tick label style={font=\footnotesize}, legend pos=south west, legend style={font=\footnotesize}]
\addplot+[mark=*, mark size=1.2pt, thick] table[x=tau, y=cobertura]{\tablaG};
\addplot+[mark=square*, mark size=1.2pt, thick] table[x=tau, y=precision]{\tablaG};
\draw[dashed] (axis cs:0.63,0) -- (axis cs:0.63,1.05);
\node[anchor=south west, font=\footnotesize] at (axis cs:0.63,0.02) {$\tau = 0{,}630$};
\legend{Cobertura, Precisión del primer candidato}
\end{axis}
\end{tikzpicture}
\caption[{Cobertura y precisión según el umbral $\tau$ del reordenador}]{Cobertura y precisión según el umbral $\tau$ del reordenador. Fuente:\allowbreak{} conjunto anotado del prototipo de calibración, 14 de septiembre de 2026. Muestra:\allowbreak{} n~=~542 ítems; $\tau$ de 0,50 a 0,73 en 24 pasos. Métrica:\allowbreak{} cobertura (ítems cuyo primer candidato tiene puntaje $\geq$ $\tau$) y precisión del primer candidato entre los cubiertos. Resultado:\allowbreak{} recalculado desde los puntajes conservados, sin volver a ejecutar el prototipo. Con $\tau$ = 0,630 se resuelven 354 de 542 ítems (65,3 \%), y 307 de ellos son correctos (precisión del 86,7 \%). El umbral se eligió y se midió sobre el mismo conjunto, lo que sesga el resultado al alza. En $\tau$ = 0,73 queda 1 ítem resuelto, así que la precisión de ese extremo descansa en muy pocos casos. El umbral del sistema construido se aplica a otro puntaje.}
\label{fig:bench-F5-2}
\end{figure}
